% Independently reconstructed fixed-depth survival theorem.
% The pair in E.33 depends on the prescribed depth and modulus.
\subsection{An independent reconstruction of fixed-depth survival}

The following results reconstruct C.2.10--C.2.12 in
\path{GT_TOWER_CHECK} and A.65--A.66 in \path{SECTIONS_OVER_Q}.
All arithmetic conditions are
made simultaneous before Chebotarev is used.  At the cohomological
step we retain the literal $c_0$ condition of A.22, as well as the
possible nonzero finite arithmetic kernel.

\begin{result}{E.30}{Proved}{Exact local images for the variable exponent}
\label{coh:variable-image}
Let $d\ge2$, let $24\mid C$, and choose $h\ge2d^2$ at least every
prime divisor of $C$.  Set
\[
 n=h!,\qquad E(Z)=Cn\prod_{r=2}^d(Z^r-1),\qquad D=Cd!n^d.
\]
For $p\mid n$, put $v=v_p(n)$, $e=v_p(d)$, and $a_p=1$ for
odd $p$, $a_2=2$.  Then the map
\[
 F_p(z)=\exp(E(z)\log z)
\]
is a bijection
\[
 1+p^v+p^{v+a_p}\mathbb Z_p
 \xrightarrow{\sim}
 F_p(1+p^v)(1+p^{W_p}\mathbb Z_p),\qquad
 W_p=v_p(C)+(d+1)v+v_p(d!)+a_p+e.
 \tag{E.30.1}
\]
It preserves congruences, agrees with $z^{E(z)}$ on positive integer
inputs, and throughout this shell
\[
 v_p(z-1)=v,\qquad v_p(E(z))=v_p(D),\qquad W_p\le2v_p(D).
 \tag{E.30.2}
\]
For $p\nmid n$, let $\omega\in\mu_{p-1}(\mathbb Z_p)$ satisfy
$\omega^r\not\equiv1\pmod p$ for $2\le r\le d$.  For a fixed
integer $e_0$ modulo $p-1$, the map
\[
 z\longmapsto \omega^{e_0}
       \exp(E(z)\log(z/\omega))
 \tag{E.30.3}
\]
is a congruence-preserving bijection from $\omega(1+p\mathbb Z_p)$
to $\omega^{e_0}(1+p\mathbb Z_p)$.  It equals the ordinary variable
power on those positive integer inputs for which $E(z)\equiv e_0$
modulo $p-1$.
\end{result}
\begin{proof}
For $p\mid n$ write $z=1+p^v u$.  Factoring the lowest power from
each factor gives an exact identity of convergent series
\[
 E(z)\log z=K u^d H(u),\qquad K=Cn\,d!\,p^{dv},
\]
\[
 H(u)=\prod_{r=2}^d
 \frac{(1+p^vu)^r-1}{rp^vu}\,
 \frac{\log(1+p^vu)}{p^vu}.
 \tag{E.30.4}
\]
We first check the useful coefficient bound
$H\in1+p^{e+1}\mathbb Z_p\langle u\rangle$.  A nonconstant
coefficient of a polynomial factor, indexed by $k\ge2$, has valuation
at least $v(k-1)-v_p(k)$, since
$\binom rk/r=\binom{r-1}{k-1}/k$.  If $p\le d$, the factorial
gives $v\ge\lfloor h/p\rfloor\ge2d$, more than the possible
denominator loss.  If $p>d$, those polynomial denominators are units
and $e=0$.  The logarithmic factor has the same bound for every
$k\ge2$.  Its minimum is at least $v$ for odd $p$ and at least
$v-1$ for $p=2$: use $v_p(k)\le k-2$ in the first case and
$v_2(k)\le k-1$ in the second.  The bounds tend to infinity with
$k$, and at two $v\ge4$.  They prove the asserted bound on $H$.

Now set $u=1+p^{a_p}w$.  The resulting series has the form
\[
 u^dH(u)=H(1)+p^{a_p+e}(\lambda w+pG(w)),
 \quad \lambda\in\mathbb Z_p^\times,
 \quad G\in\mathbb Z_p\langle w\rangle.
 \tag{E.30.5}
\]
Indeed, $v_p\binom dk\ge e-v_p(k)$, and after division by
$p^{a_p+e}$ a term of degree $k\ge2$ from $u^d$ has valuation
at least $a_p(k-1)-v_p(k)\ge1$.  After its constant term is removed,
the perturbation $u^d(H(u)-1)$ has the same extra factor $p$.
The remaining linear coefficient is $d/p^e$ modulo $p$.
Integral restricted power series are one-Lipschitz.  Consequently the
equation $\lambda w+pG(w)=t$ has exactly one solution for each
$t\in\mathbb Z_p$, by the contraction
$w\mapsto\lambda^{-1}(t-pG(w))$.  Multiplication by $K$ gives the
claimed logarithmic coset.  Exponential and logarithm are inverse
isometries on $p\mathbb Z_p$ for odd $p$ and on $4\mathbb Z_2$ at
two; all expressions lie in these domains.  This proves (E.30.1).

The same polynomial factor estimates give
$v_p(z^r-1)=v+v_p(r)$, so summation gives $v_p(E(z))=v_p(D)$.
The remaining inequality follows from
\[
 2v_p(D)-W_p
   =v_p(C)+(d-1)v+v_p(d!)-a_p-e\ge0.
\]
On this shell the polynomial, logarithm, product, and exponential all
preserve congruences.  The identity between logarithms of integer
powers gives the asserted agreement with $z^{E(z)}$.

If $p\nmid n$, then $p>h\ge2d^2$, and $E(\omega)$ is a unit.
Putting $z=\omega(1+pt)$ in (E.30.3) gives
\[
 \frac{\log(F(z)/\omega^{e_0})}{p}
       =E(\omega)t+pG_0(t),\qquad
 G_0\in\mathbb Z_p\langle t\rangle.
\]
The nonlinear logarithmic coefficient $p^{k-1}/k$ is divisible by
$p$ for every $k\ge2$; the other factor is an integral polynomial.
The same contraction proves the last bijection.  Keeping $e_0$
fixed is essential: the residue exponent is not a continuous function
of an arbitrary $p$-adic input.  Its congruence with $E(z)$ is imposed
separately on the integer inputs in E.31.
\end{proof}

\begin{result}{E.31}{Proved}{The simultaneous variable-exponent arithmetic recursion}
\label{coh:variable-arithmetic}
Fix $d\ge2$, $24\mid C$, and $J\ge2d^2$ at least every prime
divisor of $C$.  Write
\[
 n_j=(j+J)!,\quad D_j=Cd!n_j^d,\quad
 E_j(Q)=Cn_j\prod_{r=2}^d(Q^r-1).
 \tag{E.31.1}
\]
There are pairwise disjoint sequences of distinct primes
$q_j,r_j,s_j>3$ for which
\[
 A=3\prod_{j\ge0}\frac{q_j^{E_j(q_j)}}{s_j^{E_j(s_j)}},\qquad
 B=-2\prod_{j\ge0}\frac{r_j^{E_j(r_j)}}{s_j^{E_j(s_j)}}
 \tag{E.31.2}
\]
form an ordinary completed unit-equation pair.  Both coordinates have
infinite support, their quotients by $3,-2$ are global $C$-th powers,
and their nonzero valuation patterns are exactly
\[
 \begin{array}{c|c}
 2&(0,1)\\3&(1,0)\\
 q_j&(E_j(q_j),0)\\r_j&(0,E_j(r_j))\\
 s_j&(-E_j(s_j),-E_j(s_j)).
 \end{array}
 \tag{E.31.3}
\]
Every prime choice is specified by finitely many congruences and
complete splitting in one explicitly given radical field; those
conditions have an actual simultaneous Galois element.
\end{result}
\begin{proof}
Set $B(D)=(D+1)^4$ and $H_j(p)=2v_p(D_j)+j+3$.  Inductively keep
pairwise coprime coefficients
\[
 a_j=3u_j^{24},\qquad b_j=-2v_j^{24},\qquad c_j=w_j^{24},
\]
with positive $u_j,v_j,w_j$ of disjoint support away from two and three,
and impose the two invariants
\[
 a_jX^{D_j}+b_jY^{D_j}=c_jZ^{D_j}
       \text{ has a solution in }(\mathbb Z_p^\times)^3
       \text{ for every }p,
 \tag{E.31.4}
\]
\[
 v_p(a_j+b_j-c_j)\ge H_j(p)\quad(p\le B(D_j)).
 \tag{E.31.5}
\]
The initial triple $(3,-2,1)$ satisfies them with $(X,Y,Z)=(1,1,1)$.

We first prove the precise field compatibility required below.
For $24\mid n$, let $F_n=\mathbb Q(\mu_n)$.  If $(a,b,c)$ is any
current triple, each of
\[
 F_n((b/c)^{1/n}),\qquad F_n((a/c)^{1/n}),\qquad
 F_n((-a/b)^{1/n})
 \tag{E.31.6}
\]
has maximal abelian subfield $F_n$.  The radicands have squareclasses
$-2,3,6$, whose square roots lie in $F_n$.  None is an odd-prime
power in $F_n$: a valuation at two or three is $1$ or $-1$, so an
odd-prime radical has that prime degree over $\mathbb Q$; its real
radical field has nonreal conjugates and is not normal.  If any such
root lay in $F_n$, multiplication by a root of unity would put its
real root in an abelian extension, a contradiction.  Similarly no
radicand is a fourth power there.  Its fourth root, after removing
rational factors and roots of unity already in $F_n$, gives respectively
$2^{1/4}$, $3^{1/4}$, or $6^{1/4}$.  For the third radicand one uses
$(3/2)^{1/4}=6^{1/4}/\sqrt2$.  These irreducible real quartic fields
are nonnormal as well.

Kummer theory therefore gives a cyclic kernel $H$ of order $n/2$ in
the Galois group $G$ of each field (E.31.6).  For clarity, if the kernel
had order $m$, the $m$-th power of the chosen $n$-th root would belong
to $F_n$; the radicand would be an $(n/m)$-th power.  The preceding
exclusions and the existing square root force $n/m=2$.
Conjugation on $H$ has cyclotomic multiplier.  Every multiplier is odd,
and complex conjugation gives $-1$, so $[H,G]=2H$.  The fixed field of
$2H$ contains the splitting field of the relevant pure quartic, which
is a nonabelian dihedral extension of degree eight.  Hence $G/2H$ is
not abelian.  Since $G/H$ is abelian and $H/2H$ has order two, it follows
that $[G,G]=H$, proving the maximal-abelian assertion.

In particular, for $n\mid M$ the intersection of a field (E.31.6)
with $\mathbb Q(\mu_M)$ is exactly $F_n$.  Any CRT unit class equal
to one modulo $n$ defines a cyclotomic automorphism agreeing on that
intersection with the identity of (E.31.6).  The two automorphisms give
an element of the fibre-product Galois group of the compositum.
Chebotarev supplies arbitrarily large primes with both required
restrictions \cite[Theorem 8.31, p.147]{CohMilneANT}.  The restrictions
are preserved under conjugacy, so the Frobenius conjugacy-class
formulation suffices.

At stage $j$, abbreviate $n=n_j$, $D=D_j$, $E=E_j$ and
$(a,b,c)=(a_j,b_j,c_j)$.  Choose $P\ge B(D_{j+1})$ at least every
old support prime.  Before processing the primes $p\le P$ in increasing
order, fix precisions
\[
 L_p\ge H_{j+1}(p),\quad p^{L_p}>P,\quad
 L_p\ge v_p(n)+a_p\quad(p\mid n).
 \tag{E.31.7}
\]
For each of the three future prime variables we require, locally,
\[
 \begin{cases}
 z\in1+p^{v_p(n)}+p^{v_p(n)+a_p}\mathbb Z_p,&p\mid n,\\
 z\in\mathbb Z_p^\times,\quad z^r\not\equiv1\pmod p
                         \ (2\le r\le d),&p\nmid n.
 \end{cases}
 \tag{E.31.8}
\]
The second case has $p>j+J\ge2d^2$, whereas the forbidden set has
size at most $\sum_{r=2}^dr<d^2$, so permitted residues exist.

When a prime $p$ is reached, the earlier conditions already fix each
future input modulo $p-1$: every prime-power divisor of $p-1$ is
smaller than $P$, its prime was processed earlier, and (E.31.7) fixed
more than that precision.  Thus they fix separately the three classes
$e_i=E(z_i)\bmod(p-1)$.  At every prime $\ell\mid p-1$,
(E.30.2) and (E.31.8) give
$v_\ell(E(z_i))=v_\ell(D)$ for inputs consistent with the choices.
For $\ell\nmid n$ both sides are zero, since all prime divisors of
$C d!$ divide $n$.  It follows that
\[
 \gcd(e_i,p-1)=\gcd(D,p-1),\qquad
 (\mathbb F_p^\times)^{e_i}=(\mathbb F_p^\times)^D.
 \tag{E.31.9}
\]
Every forbidden residue has order dividing some $r\le d$ and $p-1$;
that order divides $\gcd(D,p-1)$, and hence every $e_i$.  All forbidden
residues therefore map to one under each power map in (E.31.9).

We now solve the local equation
$aF_a(z_a)+bF_b(z_b)=cF_c(z_c)$ without changing earlier congruences.
If $p\mid n$, use the common shell map (E.30.1).  Starting with its
common value $F_0=F_p(1+p^{v_p(n)})$, the error is
$(a+b-c)F_0$, divisible by $p^{H_j(p)}$ and therefore by $p^{W_p}$.
At least one coefficient is a unit.  Vary that output in the entire
coset $F_0+p^{W_p}\mathbb Z_p$ to cancel the error, then take its
unique preimage under E.30.  Here $p\le B(D)$, as $D\ge n$.

Suppose $p\nmid n$.  Use (E.30.3) with the separately fixed $e_i$.
If $p\le B(D)$, choose a common value
$\lambda\in(\mathbb F_p^\times)^D$ and permitted preimages for the
three maps.  If this image has more than one element, choose
$\lambda\ne1$; no forbidden preimage can occur.  If it has one
element, all nonzero residues are preimages, and the forbidden-set
bound still leaves a choice.  The three output balls are then the same
ball $\widetilde\lambda(1+p\mathbb Z_p)$.  At their common
Teichmuller value the error is divisible by $p$, by (E.31.5).  Vary an
output with unit coefficient inside this ball to cancel it exactly.
If instead $p>B(D)$, start with the unit solution (E.31.4).  Scale its
three coordinates by one common unit so that none of the three
$D$-th-power residues equals one.  There are at most three forbidden
values of the scaling power, whereas
$(p-1)/\gcd(D,p-1)>3$.  The resulting three output values have
permitted preimages by (E.31.9), and the local bijections E.30 attain
these exact $p$-adic unit values.  This also covers old coefficient
primes larger than $B(D)$, since (E.31.4) gives unit coordinates there.

Replace these inputs by residues modulo $p^{L_p}$.  Congruence
preservation retains an error divisible by $p^{H_{j+1}(p)}$.
The residues fix the exponent classes at later primes, so this
triangular process is consistent throughout.  The resulting three
classes modulo $M=\prod_{p\le P}p^{L_p}$ are units and equal one
modulo $n$.  Apply the already proved field-intersection calculation
to select distinct $q,r,s>P$ in these classes, splitting respectively
in the three fields (E.31.6).  Define
\[
 a'=a q^{E(q)},\qquad b'=b r^{E(r)},\qquad c'=c s^{E(s)}.
 \tag{E.31.10}
\]
The support conditions and twenty-fourth-power forms persist.

At $p\le P$ the new error has precision $H_{j+1}(p)$.
Vary a coordinate with unit coefficient at $(1,1,1)$.
Its derivative has valuation $v_p(D_{j+1})$, and the error has
strictly greater than twice that valuation.  The elementary strong
Hensel criterion yields a nearby unit solution for exponent $D_{j+1}$.

At a new prime $Q\in\{q,r,s\}$ the shell conditions give
$v_\ell(Q-1)=v_\ell(n)$ for every $\ell\mid n$.  A prime divisor
of $D_{j+1}$ not dividing $n$ is at most $j+J+1\le P$; its second
condition in (E.31.8), already with $r=2$, excludes $Q\equiv1$
modulo that prime.  Since $n\mid D_{j+1}$, we obtain the exact identity
\[
                  \gcd(D_{j+1},Q-1)=n.
 \tag{E.31.11}
\]
Thus the $D_{j+1}$-power and $n$-th-power images in $\mathbb F_Q^\times$
coincide.  At $q$, the split radical makes $c/b$ an $n$-th power, and
the newly introduced $r,s$ powers are also $n$-th powers.  With
$X=Z=1$, solve $Y^{D_{j+1}}=(c'-a')/b'$ by simple-root Hensel.
At $r$, take $Y=Z=1$ and use the split ratio $c/a$.
At $s$, take $Y=Z=1$ and solve
$X^{D_{j+1}}=(c'-b')/a'$ using the split ratio $-b/a$.
All roots and derivatives are units because the new primes exceed
$B(D_{j+1})$.

At any remaining prime $p>B(D_{j+1})$ all coefficients are units and
$p\nmid D_{j+1}$.  The smooth projective diagonal curve of degree
$D_{j+1}$ has genus $(D_{j+1}-1)(D_{j+1}-2)/2$.  The Weil bound
\cite[(1.5.1) and Theorem (1.6), pp.275--276]{RootDeligne1974}
and at most $3D_{j+1}$ boundary points give at least
\[
 p+1-(D_{j+1}-1)(D_{j+1}-2)\sqrt p-3D_{j+1}>0
\]
points with all coordinates nonzero.  To check the strict inequality,
put $t=\sqrt p>(D_{j+1}+1)^2$ and replace
$(D_{j+1}-1)(D_{j+1}-2)$ by $D_{j+1}^2$; then
$t^2-D_{j+1}^2t-3D_{j+1}+1>
t(2D_{j+1}+1)-3D_{j+1}+1>0$.  A nonzero coordinate derivative
lifts such a point to a unit solution.  This proves both inductive
invariants at every prime.  The curve bound is an independent
replacement for the Jacobi-sum estimate in the historical proof.

Finally every new exponent is divisible by $Cn_j$.  In the global
power completion, the products and their versions with exponents
divided by $C$ converge: modulo any fixed positive integer power,
every sufficiently late factor is trivial.  At a fixed finite prime
$p$ each coefficient valuation eventually stabilizes.  Later factors
are units, and growing factorial exponents annihilate every finite
quotient of $\mathbb Z_p^\times$.  Hence the coefficient sequences
converge in $\mathbb Q_p^\times$ to nonzero $a_p,b_p,c_p$.
Condition (E.31.5) applies eventually to this prime with precision
tending to infinity, and gives $a_p+b_p=c_p$.  The ordinary point
$t_p=a_p/c_p$ has $1-t_p=b_p/c_p$ and precisely the localized
completed pair (E.31.2).  At the real place the signs are $(+,-)$,
represented for example by $(3,-2)$.  Distinctness of the support
primes proves (E.31.3), and the divided convergent products prove the
$C$-th-power assertion.  This completes the recursion and the audit
of its continuity and field compatibility.
\end{proof}

\begin{result}{E.32}{Proved}{Uniform Hall, inertia, and arithmetic-kernel bounds}
\label{coh:fixed-depth-bounds}
Fix $R\ge3$.  Let $N_r=\Pi/\Gamma_{r+1}\Pi$ and
$L_r=\Gamma_r\Pi/\Gamma_{r+1}\Pi=\widehat{\mathbb Z}(r)^{b_r}$.
Use the actual point-three section $s_3$ to identify
$E_r=N_r\rtimes_{\rho_3}G_{\mathbb Q}$.  Define
\[
 K_r(m)=\ker\left(N_r\longrightarrow
 (1+I)((\mathbb Z/m)\langle X,Y\rangle/I^{r+1})\right),
 \qquad I=(X,Y),
 \tag{E.32.1}
\]
using the integral Magnus images $x\mapsto1+X$, $y\mapsto1+Y$.
If $h_R$ is the number of Hall basis elements through weight $R$, put
\[
 D_H=(R!)^{h_R+1},\quad T=2(R!)^R,\quad
 B_r=2|K_{2r-2}(\mathbb Z)|\quad(2\le r\le R).
 \tag{E.32.2}
\]
Then $K_r(m)$ is characteristic, the continuous Hall set section
$\sigma_r:N_{r-1}\to N_r$ with zero new coordinates satisfies
\[
 \sigma_r(K_{r-1}(D_Hm))\subset K_r(m),
 \tag{E.32.3}
\]
and the finite integer $B_r$ annihilates
$H^2(G_{\mathbb Q,\{\ell,\infty\}},\mathbb Z_\ell(r))$ for every
coefficient prime $\ell$.

At a fresh prime $q>3$, $q\ne\ell$, suppose an ordinary local
abelian pair has valuation pattern $(E,0)$, $(0,E)$, or $(-E,-E)$.
For any actual class-$(r-1)$ section with this pair, the central
defect $c$ of its canonical Hall set lift satisfies the strict
tame-residue bound
\[
 \frac{z_q(c)}{q^{r-1}-1} \in
 \frac{E}{T\prod_{k=2}^{r-1}(q^k-1)}
                      \mathbb Z_\ell^{b_r}.
 \tag{E.32.4}
\]
The product is empty for $r=2$.  We use this assertion when all the
displayed ideals are integral.  At a prime $q>3$, $q\ne\ell$,
where both ordinary valuations are zero, the defect residue is zero.
\end{result}
\begin{proof}
The Hall model and its pro-$\ell$ universal property were established
in A.20.  The product over $\ell$ gives $N_r$: a finite nilpotent
group is the product of its Sylow groups, and these are precisely its
finite nilpotent quotients of class at most $r$.  On a weight-$r$
commutator the actual Galois action is $\chi^r$; conjugating the
tangential action by the point-three section does not change this
central graded action.

Substituting the integral Magnus series of any two generator images
preserves the ideal of coefficients divisible by $m$.  This proves
that $K_r(m)$ is characteristic.  To check (E.32.3) without excluding
small coefficient primes, use a Lyndon Hall basis.  Its lowest Magnus
polynomials have distinct leading words with coefficient one; successive
Hall exponent extraction is integral and triangular.  If an exponent
$a$ is divisible by an integral ideal $J$, every nonconstant coefficient
of its Hall power through degree $R$ is divisible by $J/(R!)$:
\[
           \binom ak=\frac ak\binom{a-1}{k-1},\qquad k\le R.
\]
The identity applies in $\mathbb Z_\ell$, including when $\ell\mid k$.
Extract an exponent and remove its Hall power successively.  There are
at most $h_R$ removals, each losing at most $R!$ in this conservative
bound.  Starting with coefficients divisible by $D_Hm$, every
extracted exponent is still divisible enough that its extension
through degree $r$ is one modulo $m$.  Multiplying the extended
Hall powers proves (E.32.3).

Quillen's finite-generation theorem and Borel's rank formula show
that each even positive group $K_{2r-2}(\mathbb Z)$ is finite
\cite[IV.6.9, author p.59; IV.1.17--1.18, author p.12]{CohWeibelKbook}.
At odd $\ell$, the integral comparison identifies the displayed
$H^2$ with its $\ell$-primary part
\cite[VI.8.2, author p.47]{CohWeibelKbook}.
Inverting $\ell$ does not change that primary part: localization adds
only the adjacent finite-field $K$-groups, whose nonzero orders are
$\ell^i-1$ \cite[IV.1.13; V.6.8]{CohWeibelKbook}.
At two, VI.9.11 of the same source gives the comparison with the
real-place term retained.  For $r\not\equiv0\pmod4$ the group is
the two-primary $K$-group (the possible integer $\rho$ is zero since
the real rank is one).  For $r\equiv0\pmod4$, its kernel under
restriction to the real place is that group and the remaining image
has exponent at most two.  Thus the additional factor two in $B_r$
annihilates the entire group in every case.  This is the precise
finite-kernel calculation used in A.22 and B.25; it neither assumes
an unknown even $K$-group to vanish nor discards a derived-limit class.

We next check that the action used in the tame bound is the actual
unramified action.  For $q\ne\ell$ the tangential action in class
two is diagonal cyclotomic: its actual path commutator lies in
$\Gamma_2$ and is central in this quotient.  Suppose the action is
unramified in class $i$.  The kernel of the next automorphism map is
a submodule of $\operatorname{Hom}(N_1,L_{i+1})$, of Frobenius
weight $i$.  Wild inertia and non-$\ell$ tame inertia map trivially
to this pro-$\ell$ torsion-free module.  For the remaining generator,
$FTF^{-1}=T^q$ makes its image annihilated by $q^i-q$.  For $i\ge2$
this integer is nonzero, so the image vanishes.  The tangential action
is therefore unramified at every class.  For $q>3$ the point-three
section is itself unramified: its abelian inertia values are zero,
and successively the same tame relation kills a first possible
nonzero coordinate in every weight at least two.  Hence $\rho_3$
is unramified at all fresh support primes.

In local class two, write the actual tame inertia as
$V_2=y^n x^m [x,y]^t$.  The explicit ordinary local cup primitives in
A.25 show
\[
                m,n\in E\mathbb Z_\ell,\qquad
                t\in(E/2)\mathbb Z_\ell.
 \tag{E.32.5}
\]
For a positive pattern the central coordinate is the valuation times
an integral primitive.  For the negative pattern it is such a product
plus a binomial in the valuation, whose denominator is at most two.
This bound holds for every local integral lift: the difference of
two central primitives is a $\mathbb Z_\ell(2)$-cocycle, and its
inertia value is killed by $q-1$, hence is zero.  Since $s_3(T)=1$,
these are also the relative inertia coordinates.  Extending $V_2$
canonically through degree $R$ is consequently one modulo
$E/(2R!)$ in the Magnus algebra.

Suppose an actual class-$i$ section is given.  Let $\bar V$ be the
canonical lift of its class-$(i-1)$ inertia and write its last
coordinate as $Z_i\in L_i$.  If $U$ is its actual Frobenius coordinate
and $\varphi_q$ the actual unramified action, the tame relation reads
\[
 U\varphi_q(\bar V)U^{-1}\bar V^{-q}
                     = Z_i^{-q(q^{i-1}-1)}.
 \tag{E.32.6}
\]
Integral substitution, multiplication, conjugation and powers preserve
any integral Magnus coefficient ideal.  The left side is central;
its weight-$i$ Hall coordinates are therefore extracted integrally,
with no denominator.  Divide by $q^{i-1}-1$ (and the unit $q$), and
allow a further factor $R!$ to extend the resulting Hall power through
degree $R$.  Starting from (E.32.5), induction gives the bound
\[
 V_i^{\rm can}\equiv1\pmod{
 E/(2(R!)^{i-1}\prod_{k=2}^{i-1}(q^k-1))}
                  \quad(2\le i\le R)
 \tag{E.32.7}
\]
for the canonical extension of actual inertia.

The central relator of the canonical lift is exactly $qz_q(c)$ by
the cochain-to-purity calculation B.22.  Use (E.32.7) at $i=r-1$,
and divide once more by $q^{r-1}-1$; the weaker denominator $T$
in (E.32.4) absorbs all factorial losses.  For $r=2$, an actual local
class-two lift exists by A.25 and (E.32.5) supplies its correction,
giving the stronger bound $z_q(c)/(q-1)\in(E/2)\mathbb Z_\ell$.
If both valuations are zero, all successive inertia coordinates vanish
by the torsion-free tame relation, so the canonical relator is zero.

All cochains here lie in the actual tame domain of B.22.  Indeed the
tangential action is unramified, and a local pro-$\ell$ section kills
wild and prime-to-$\ell$ tame inertia.  Its relative cocycle and its
Hall set lift factor through that same tame quotient.  This observation
also applies at $q=2,3$ when $q\ne\ell$, although $\rho_3$ there may
have inner inertia.  In particular the residue bounds are bounds on
the strict global localization representative, not just on a formal
relator in a diagonally acted group.
\end{proof}

\begin{result}{E.33}{Proved}{Ordinary infinite support at every prescribed integral depth}
\label{coh:fixed-depth}
For every integer $R\ge1$ and positive integer $N$ there is an ordinary
infinite-support completed pair $(A,B)$ and a continuous section of
the actual integral class-$R$ quotient with that abelianization.
It agrees with the point-three section in the finite arithmetic quotient
defined by (E.32.1) modulo $N$.  Its pair is neither a rational-point
pair nor a cuspidal pair.  The pair and the section may depend on $R,N$.
This proves the fixed-depth content of GT\_TOWER\_CHECK C.2.12 and
SECTIONS\_OVER\_Q A.65, with C.2.10--11/A.66 reconstructed above.
\end{result}
\begin{proof}
It suffices to treat $R\ge3$, and project for smaller depth.  Let
$D_H,T,B_r$ be as in E.32 and put
\[
 D_r^{\rm old}=(2^{r-1}-1)(3^{r-1}-1),\qquad
 M_R=N,\qquad M_{r-1}=D_H B_r M_rD_r^{\rm old}\quad(2\le r\le R).
 \tag{E.33.1}
\]
Choose $C$ divisible by $24$, $M_1$, and all $TB_rM_r$.
Apply E.31 with $d=R-1$ and an allowed $J$.  At each fresh support
prime $q$ its absolute nonzero valuation is
\[
 E_q=C e_q\prod_{k=2}^{R-1}(q^k-1),\qquad
 e_q=n_j\ \text{on the three primes of stage }j,\qquad
 e_q\longrightarrow0\text{ in }\widehat{\mathbb Z}.
 \tag{E.33.2}
\]
The relative classes $A/3,B/(-2)$ are global $M_1$-th powers.  Choose
literal relative Kummer cocycles pointwise divisible by $M_1$;
these start the class-one section induction.

Suppose a section through class $r-1$ has relative cocycle $z$ with
$z(g)\in K_{r-1}(M_{r-1})$ for every $g$.  Lift it by $\sigma_r$.
By (E.32.3) its values lie in $K_r(L)$, where
$L=B_rM_rD_r^{\rm old}$.  Its normalized central multiplication
defect $c$ has values in $LL_r$: all Magnus coefficients of the defect
are divisible by $L$, and central Hall extraction is integral.
Thus $c/(B_rM_r)$ is an actual integral continuous two-cocycle.

Fix a coefficient prime $\ell$.  At an old support prime $q=2,3$,
$q\ne\ell$, the cochain formula B.22 makes the strict residue of
$c/(B_rM_r)$ divisible by $D_r^{\rm old}$.  Division by the needed
$q^{r-1}-1$ is therefore integral; the factor $q^{-1}$ in that formula
is an $\ell$-adic unit.  At a fresh support prime $q\ne\ell$, E.32
and (E.33.2) give
\[
 \frac{z_q(c/(B_rM_r))}{q^{r-1}-1}
 \in
 \frac{C e_q\prod_{k=r}^{R-1}(q^k-1)}{TB_rM_r}
                   \mathbb Z_\ell^{b_r}
 \ \subset\ e_q\mathbb Z_\ell^{b_r}.
 \tag{E.33.3}
\]
At every other prime away from $\ell$ the residue is zero.  The
coordinatewise quotients are integral and tend to zero, hence belong
to the literal $c_0$ module of A.22.  Its exact localization sequence
therefore places
\[
                         \beta=[c/(B_rM_r)]
 \tag{E.33.4}
\]
in the image of $H^2(G_{\mathbb Q,\{\ell,\infty\}},
\mathbb Z_\ell(r)^{b_r})$.  E.32 gives $B_r\beta=0$.
It follows that $[c]=0$, so there is a continuous integral cochain
$h$ with $\delta h=-c$ and hence an actual class-$r$ lift.

We must retain the finite congruence, rather than merely produce some
lift.  Reduction of $h$ modulo $B_rM_r$ is a cocycle.  The two
coefficient connecting maps satisfy the literal identities
\[
 \partial_{B_rM_r}[\bar h]=-\beta,\qquad
 \partial_{M_r}[\bar h\bmod M_r]=-B_r\beta=0.
 \tag{E.33.5}
\]
The second identity is just
$\delta h/M_r=B_r(\delta h/(B_rM_r))$.  Exactness for
$0\to L_r\xrightarrow{M_r}L_r\to L_r/M_r\to0$ now lifts the
reduced cohomology class to an integral one-cocycle $z_r$.  If the
representatives differ by a finite coboundary, lift its zero-cochain
and adjust $z_r$ by the corresponding integral coboundary.  Thus
$z_r\equiv h\pmod{M_r}$ pointwise.  Replacing $h$ by $h-z_r$
preserves $\delta h=-c$ and makes every new central coordinate
divisible by $M_r$.  Together with the original set lift this gives
a class-$r$ cocycle pointwise in $K_r(M_r)$, completing the induction.

These exact sequences are sequences of continuous cochain complexes:
a finite-valued cochain lifts by choosing lifts of its finitely many
values.  The argument can be carried out in every $\ell$-component.
A map into a product is continuous exactly when its components are
continuous, so the corrected component cocycles form one continuous
section of the full integral quotient.  This also justifies the full
$\widehat{\mathbb Z}$ formulation of (E.33.5).

For precision, the actual finite arithmetic quotient is
\[
 E_R\longrightarrow
 (N_R/K_R(N))\rtimes\operatorname{im}(\bar\rho_3),\qquad
 (z,g)\longmapsto(z\bmod K_R(N),\bar\rho_3(g)),
 \tag{E.33.6}
\]
where $\bar\rho_3$ is the induced finite action.  The characteristic
property in E.32 makes this a well-defined homomorphism.  Since
$M_R=N$, its two section maps are literally equal, as claimed.

Finally E.31 gives infinitely many nonzero valuation coordinates,
so the pair cannot come from a rational point.  The actual abelian
cusp map has second coordinate trivial at zero, first coordinate
trivial at one, and equal valuation coordinates at infinity.
The permanent patterns $(0,1)$ at two and $(1,0)$ at three exclude
each case.  Ordinary localization is supplied by E.31, independently
of any assertion that the particular higher section is locally a
point section.  No such higher local matching is inferred.  In
particular this theorem does not choose one pair for all $R$; applying
the fixed-pair compactness theorem A.21 to this varying family would
change its quantifiers and is invalid.
\end{proof}

\begin{result}{E.34}{Proved}{Why the denominator-clearing recursion does not diagonalize}
\label{coh:diagonal-barrier}
For a fixed prime $q$, put
\[
 P_R(q)=\prod_{k=2}^{R-1}(q^k-1)\qquad(R\ge3).
\]
For every coefficient prime $\ell\ne q$,
$v_\ell(P_R(q))\to\infty$.  Consequently
\[
 \bigcap_{R\ge3}P_R(q)\mathbb Z=\{0\},\qquad
 \bigcap_{R\ge3}P_R(q)\widehat{\mathbb Z}
 =\{z\in\widehat{\mathbb Z}:z_\ell=0\text{ for all }\ell\ne q\}.
 \tag{E.34.1}
\]
Thus the sufficient exponent divisibilities imposed in E.33 cannot be
maintained at a fixed nonzero ordinary valuation through unbounded
nilpotent depth. Moreover, the abelian outputs of the displayed
E.33 construction converge in the global power completion to
$(3,-2)$ as $R\to\infty$, for arbitrary choices of its finite
moduli $N_R$. This also verifies the exact convergence assertion
GT C.2.1 about the displayed SECTIONS A.65 parameters.
Neither statement says that the exponent divisibilities are necessary
for actual nilpotent survival.
\end{result}
\begin{proof}
For odd $\ell\ne q$, let $t$ be the multiplicative order of $q$
modulo $\ell$.  Every positive multiple of $t$ contributes at least
one to the valuation of $P_R(q)$, except that the omitted index $k=1$
must be removed if $t=1$.  Hence
\[
 v_\ell(P_R(q))\ge
 \left\lfloor\frac{R-1}{t}\right\rfloor-\mathbf1_{t=1}
 \longrightarrow\infty.
\]
This also covers $q=2$, since then every other coefficient prime is
odd.  If $\ell=2\ne q$, the prime $q$ is odd and $q^2\equiv1\pmod8$.
Every even index contributes at least three, giving
\[
 v_2(P_R(q))\ge3\left\lfloor\frac{R-1}{2}\right\rfloor
 \longrightarrow\infty.
\]
There is no factor divisible by $q$ itself.  These observations prove
the componentwise intersection in (E.34.1).  Its intersection with
the diagonal copy of $\mathbb Z$ is zero, since a nonzero integer has
finite valuation at every prime.  In particular, an ordinary support
coordinate $E_q\ne0$ cannot satisfy $P_R(q)\mid E_q$ for all $R$.
The actual sufficient condition (E.33.2) is still stronger, so has
this obstruction already before its other congruence moduli are used.

For the exact historical convergence assertion, the constants in
(E.32.2) and (E.33.1) satisfy
$R!\mid D_H=(R!)^{h_R+1}\mid M_1\mid C_R$.
The divided products in E.31 show that $A_R/3$ and
$B_R/(-2)$ are global $C_R$-th powers.
Given any positive integer $m$, every $R\geq\max(3,m)$
has $m\mid R!\mid C_R$. Their images in
$\mathbb Q^\times/(\mathbb Q^\times)^m$ are therefore
trivial for all sufficiently large $R$, independently
of the moduli $N_R$.
This is exactly convergence to $(3,-2)$ in the
inverse-limit power topology. Infinite support at each
finite stage is compatible with this rational limiting pair.

The argument concerns the mechanism that clears all displayed tame
denominators by divisibility of the valuation.  It does not apply to
a different mechanism in which the local central numerator itself
cancels those denominators.  Actual local point sections provide such
cancellation; A.31 and B.25 explain its integral and continuous forms.
Therefore this obstruction neither refutes an ordinary all-depth ghost
nor proves finite support for such a ghost.  It pinpoints the missing
step in an attempted diagonalization of E.33 or of its historical
fixed-depth predecessors.
\end{proof}
